% ============================================================================
% EVALUATION — restructured draft (RESISC45 / 8GB-satellite model)
% Night-run draft. Numbers from data/jetson_resisc45/ + regenerated sim CSVs.
% Markers [Q6?] flag reframings that depend on OPEN_QUESTIONS Q6 (user's call).
% Not yet merged into main.tex / Overleaf.
% ============================================================================
\section{Performance Evaluation}\label{EVAL}

We evaluate the framework in three signposted parts so the reader can place each
figure. We first \emph{compare} the proposed scheduler against four baselines on the
constellation, isolating where it leads and where it only trades off. We then
\emph{validate} the theoretical guarantees against a certified offline optimum and the
stability conditions. We finally \emph{stress-test} robustness under intermittent
inter-satellite links and varying eclipse. A comparison figure always shows all four
baselines; a validation or stress-test figure reports the proposed scheduler against the
relevant theoretical reference.

\subsection{Hardware Calibration and the Quality--Energy Landscape}\label{EV_CALIB}
We profile per-configuration inference cost on an NVIDIA Jetson Orin NX (16\,GB, JetPack~6,
\texttt{llama.cpp} CUDA backend, \texttt{Q4\_K\_M} weights). The task is 45-class
remote-sensing scene classification on RESISC45~\cite{Cheng2017RESISC45}. Latency is timed
per image; power is sampled every 50\,ms with \texttt{tegrastats}; energy integrates power
over the inference window; the peak bound $\hat{P}_{im}^{\text{CP}}$ is the 95th percentile.

We profiled seven vision-language models spanning four families and 2--8B parameters. Two
findings drive the rest of the evaluation. First, accuracy is \emph{not} monotone in model
size: Qwen2-VL-2B (0.398) exceeds Qwen2.5-VL-3B at the EuroSAT scale and several 4B models,
because architecture and model generation matter as much as parameter count
(Fig.~\ref{fig:tradeoff}). A configuration's quality therefore cannot be predicted from its
size and must be measured; the scheduler reasons over the measured operating points, and its
action set is the Pareto-efficient subset of profiled configurations
(Table~\ref{tab_calib}). Second, the deployment target is not the 16\,GB profiling host but a
typical sensing satellite with an $\sim$8\,GB edge module whose memory is shared with the
on-board detection pipeline and I/O, leaving only a few gigabytes for inference. Under this
budget any configuration of 4B parameters or more cannot run standalone and must be split
across two satellites; the highest-quality reachable model, Qwen2.5-VL-7B, is exactly such a
split-only configuration. This is what makes the two-stage pipeline necessary rather than
optional.

\begin{table}[!t]
\centering\small
\caption{\label{tab_calib}Measured per-configuration cost on Jetson Orin NX (RESISC45,
\texttt{Q4\_K\_M}). The scheduler's action set is the Pareto-efficient subset; the 7B tier is
split-only under the 8\,GB satellite budget. $N_{im}$ is the number of $\tau{=}1$\,s slots.}
\begin{tabular}{lccccc}
\toprule
Config & $Q_{im}$ & $\mathbb{E}[T_{im}]$(s) & $\mathbb{E}[E_{im}]$(J) & $\hat{P}_{im}^{\text{CP}}$(W) & $N_{im}$ \\
\midrule
Qwen2-VL-2B   & 0.398 & 1.24 & 3.93 & 4.15 & 2 \\
Qwen2.5-VL-3B & 0.464 & 1.74 & 5.78 & 4.51 & 2 \\
Qwen2.5-VL-7B & 0.573 & 2.76 & 9.13 & 5.57 & 3 \\
\bottomrule
\end{tabular}
\end{table}

\subsection{Benchmark Solutions}\label{EV_BASE}
We compare the proposed scheduler against four online baselines, each relaxing one or more of
the four structural constraints so that the comparison isolates the value of respecting them.
Greedy-Q always selects the highest-quality standalone configuration; it ignores energy
sustainability, so under tight resources it drains the eclipse-side battery, and because it
omits the peak-power check it can exceed the instantaneous cap. Greedy-E always selects the
cheapest configuration; it stays energy-safe but never reaches the higher-quality tiers.
Random selects a feasible standalone configuration uniformly at random, an uninformed lower
bound. Static routes to the highest-battery satellite and runs a single fixed standalone
configuration, the natural non-adaptive heuristic; because it never splits, it is capped
below the top quality tier. None of the four can orchestrate the two-satellite split, so all
are limited to standalone configurations and cannot run the 7B tier.

We treat the offline MILP solely as a clairvoyant upper bound for the optimality-gap study
(Section~\ref{EV_GAP}), not as a competing policy: it is non-causal, requiring full future
knowledge of solar arrivals, channels, and eclipse, so comparing it on equal footing with
causal online schedulers would be unfair. It is moreover intractable at constellation scale,
which is why we solve it only on a small certified-optimal instance.

% --------------------------------------------------------------------------
% §D Simulation Results — drafted with NEW RESISC45 numbers. [Q6?] = reframing
% pending user confirmation (OPEN_QUESTIONS Q6).
% --------------------------------------------------------------------------
\subsection{Multi-Satellite Comparison}\label{EV_MULTI}
On the four-satellite constellation with the split pipeline (Table~\ref{tab_split},
Fig.~\ref{fig:split}) the proposed scheduler is the only policy that runs the 7B tier,
executing 3755 inferences as a two-satellite split; the four baselines are structurally
capped at the 3B standalone tier. % [Q6?] At this default operating point the RESISC45
% configurations are inexpensive, so no policy blacks out here; the sustainability advantage
% is therefore reported where it binds, under tightened resources, in Section~\ref{EV_SENS}.
Cumulative quality is an honest trade-off rather than a blanket win: energy-aware baselines
reach higher raw throughput by running a cheap fixed configuration at full duty, whereas the
proposed scheduler spends part of its budget on the expensive 7B split and optimizes the
max-min objective. % numbers: ours totQ=3388, split=3755; Static/Greedy-Q totQ=3703.

% [Q6?] Peak-power shield: present as a by-construction guarantee that binds under a tight
% cap (Section~\ref{EV_SENS} shows baselines violating ~16k times once the cap drops below
% the 5.57W draw of the configs), rather than a headline empirical win, since the RESISC45
% peak draws are low and close together. CONFIRM framing with user (Q6).

% TODO (post-Q6): EV_SINGLE (Fig4 + Greedy-Q overlay), EV_GAP (Fig6, absorb Fig5),
% EV_SENS (Fig9/10, 2x4), EV_ISL (Fig8), EV_STAB (Fig11/12). Fill numbers from regenerated
% sim CSVs once figures are restyled.
